\documentclass[11pt]{article}
\usepackage[margin=1in]{geometry}
\usepackage{amsmath,amssymb,bm}
\title{Electric potential in the cavity and outside}
\author{}
\date{}
\begin{document}
\maketitle

\section{Problem Setup}
The conductor potential $V_c$ is already known from the conduction
calculation. The same mesh also covers the air inside the loop's cavity
and the air outside the conductor. These are separate regions; together
call them $\Omega_a$. The app assumes uniform permittivity $\epsilon$
and no volume charge in air. Gauss's law and
$\bm E=-\nabla V$ give
\begin{equation}
 \nabla\cdot(\epsilon\bm E)=0
 \quad\Longrightarrow\quad
 \nabla^2V=0\quad\text{in }\Omega_a.
\end{equation}
The distant rectangular boundary is held at $V=0$ to approximate open
space. Potential is continuous at a conductor--air surface, so its air
value there must equal the conductor value.

There is one ambiguity: conduction determines $V_c$ only up to a
constant. Shifting the conductor potential relative to the fixed far
boundary changes the electric field outside. Write the conductor--air
boundary value as $V_c+C$. The unknown $C$ will be chosen so the
circuit has zero net charge.

\section{Solving the air equation for a given boundary value}
For air, the generic equation in \texttt{fem.tex} has $c=1$ and
$\bm f=0$. Every air boundary has a prescribed potential: the
conductor surfaces have $V_c+C$, and the far rectangle has zero.
We choose test functions $w$ as the basis functions of interior air vertices.
These basis functions are zero on the boundary, so the boundary integral
term in \texttt{fem.txt} vanishes.
The resulting free-vertex equations are
\begin{equation}
 \int_{\Omega_a}\nabla V\cdot\nabla w\,dA=0.
\end{equation}
Each air triangle contributes
$K^{(h)}_{ij}=A_h\nabla N_{h,i}\cdot\nabla N_{h,j}$, with no volume-source
contribution. The same assembled system contains both the cavity
and outside triangles. Each air region has its own prescribed
boundary values.

We do not yet know $C$, but this equation is linear. Solve it first
for two sets of boundary values:
\begin{align}
 U|_{\rm conductor}&=V_c,& U|_{\rm far}&=0,\\
 H|_{\rm conductor}&=1,& H|_{\rm far}&=0.
\end{align}
Both $U$ and $H$ satisfy the zero-source air equation. The solution
for any offset is therefore
\begin{equation}
 V_a=U+CH.
 \label{eq:combination}
\end{equation}
Inside the enclosed cavity $H=1$, so changing $C$ shifts the cavity
potential without changing its electric field. Outside, $H$ changes
from one at the conductor to zero at the far boundary.

\section{Finding the offset from charge neutrality}
The far boundary surrounds the whole circuit. By Gauss's law, the
total enclosed charge per unit depth is
\begin{equation}
 Q'=\epsilon\oint_{\Gamma_{\rm far}}\bm E\cdot\bm n\,ds
   =-\epsilon\oint_{\Gamma_{\rm far}}\nabla V\cdot\bm n\,ds.
\end{equation}
We want $Q'=0$, so we need the total outward flux of $\nabla V$
through $\Gamma_{\rm far}$. To obtain it from the air solution,
let $\mathcal F$ be the set of far-boundary vertices and choose
\begin{equation}
 W_{\rm far}=\sum_{i\in\mathcal F}N_i.
\end{equation}
The basis functions sum to one along the far boundary, so
$W_{\rm far}=1$ there. They are zero along the conductor surfaces,
so $W_{\rm far}=0$ there. The integration-by-parts equation in
\texttt{fem.tex}, with $c=1$ and $\bm f=0$, now gives
\begin{align}
 \int_{\Omega_a}\nabla V\cdot\nabla W_{\rm far}\,dA
 &=\int_{\partial\Omega_a}W_{\rm far}\nabla V\cdot\bm n\,ds\\
 &=\oint_{\Gamma_{\rm far}}\nabla V\cdot\bm n\,ds.
\end{align}
The cavity contributes no boundary term because $W_{\rm far}=0$
on its conductor boundary. Thus the total charge is $-\epsilon$
times this volume integral.

For the computed piecewise-linear potential, use the triangle
expansion from \texttt{fem.tex}:
\begin{equation}
 V_h=\sum_{k=0}^2V_{h,k}N_{h,k}.
\end{equation}
Writing $W_{{\rm far},h}$ for the restriction of $W_{\rm far}$
to triangle $h$, the corresponding flux is
\begin{align}
 \Phi(V)
 &:=\int_{\Omega_a}\nabla V\cdot\nabla W_{\rm far}\,dA\\
 &=\sum_h\sum_{k=0}^2
   A_hV_{h,k}\nabla N_{h,k}\cdot\nabla W_{{\rm far},h}.
 \label{eq:reaction}
\end{align}
This is linear in the potential values. The far-boundary
basis functions were excluded as test functions when solving for
those values, but we can use their sum to calculate the flux afterward.

Using equation~\eqref{eq:combination}, this flux is
$\Phi(U+CH)=\Phi(U)+C\Phi(H)$. Set it to zero and solve the one remaining
scalar equation:
\begin{equation}
 C=-\frac{\Phi(U)}{\Phi(H)}.
\end{equation}
The uniform factor $\epsilon$ cancels. This is why the code solves
the air system twice: one solve gives the field caused by the
conductor's varying potential, and the other gives the response to
a uniform potential shift.

\section{Completing the electric field}
Add $C$ to the conductor's potentials and assign $U+CH$ to
air vertices. Conductor and air share boundary vertices, so the
combined potential is continuous. For each air triangle,
\begin{equation}
 \bm E_h=-\sum_{i=0}^2 V_{h,i}\nabla N_{h,i}.
\end{equation}
The constant shift leaves the conductor electric field unchanged.
The app uses the completed field both inside and outside the circuit.

\paragraph{Code:} \texttt{calculations/exterior.py}.
\end{document}
